\subsection{双向 KL 常数、Gödel 外置的输运障碍与有限 CMV 证书的 Pad\'e 冻结}\label{subsec:conclusion-godel-transport-cmv-freezing-complexity-gap}
在残差外置的条件 Kolmogorov 复杂度、边界均匀口径下的几何熵率、Gödel 单整数化、面积保持椭圆输运以及有限 CMV 证书塔的点态稳定化都已经建立之后，还可把这些链条进一步压缩为一组终端命题：双向 KL 线性常数与最优残差码率之差精确一致，中心规范熵在全纤维对称中只占零比例；Gödel 单整数化相对于最优残差账本存在不可消去的对数发散鸿沟，而这一鸿沟在椭圆输运条件数中被逐字保留。另一方面，代数尺度下的椭圆输运至多删去若干个 \(C_2\) 的 Jordan--Hölder 因子，因而既不能洗去非阿贝尔奇数部分，也不能消除非可解性。最后，点态 CMV 稳定化强迫任意固定阶 Pad\'e--Hardy 局部对象在某一有限级别后完全冻结，从而把全部局部有理谱谓词压缩为有限级别即可离线判定的命题。

\begin{theorem}[双 KL 常数与残差条件复杂度密度的同一律]\label{thm:conclusion-double-kl-residual-conditional-complexity-identity}
\leanverified{paper\_conclusion\_double\_kl\_residual\_conditional\_complexity\_identity}
设
\[
R_m(\omega):=K\!\bigl(r_{m,S}(\omega)\mid X\bigr),
\]
并假定线性增长常数 \(c_{u\to w},c_{w\to u}\) 与边界熵率 \(h_\partial\) 存在。则有
\[
\lim_{m\to\infty}\frac1m\,\EE_\omega R_m(\omega)
=
\frac{\log 2-h_\partial+c_{w\to u}}{\log 2},
\]
而边界均匀口径下的理想枚举残差码率满足
\[
\lim_{m\to\infty}\frac{1}{m\log 2}\sum_{x\in X_m}u_m(x)\log d_m(x)
=
\frac{\log 2-h_\partial-c_{u\to w}}{\log 2}.
\]
从而二者之差精确为
\[
\lim_{m\to\infty}\left[
\frac1m\EE_\omega R_m(\omega)
-
\frac{1}{m\log 2}\sum_{x\in X_m}u_m(x)\log d_m(x)
\right]
=
\frac{c_{u\to w}+c_{w\to u}}{\log 2}.
\]
\end{theorem}
\begin{proof}
正文已给出
\[
\EE_\omega R_m(\omega)=\frac{\kappa_m(\pi_m)}{\log 2}+O(1),
\]
而 \(\kappa_m(\pi_m)=m\,\gamma_m^{\mathrm{Sh}}\)。故
\[
\lim_{m\to\infty}\frac1m\EE_\omega R_m(\omega)
=
\frac{\lim_{m\to\infty}\gamma_m^{\mathrm{Sh}}}{\log 2}.
\]
再由正文中的双向 KL 常数闭式，
\[
\lim_{m\to\infty}\gamma_m^{\mathrm{Sh}}
=
\log 2-h_\partial+c_{w\to u},
\]
从而得到第一式。

另一方面，
\[
\frac{1}{m\log 2}\sum_{x\in X_m}u_m(x)\log d_m(x)
=
\frac{\gamma_m^{\mathrm{geo}}}{\log 2},
\]
而
\[
\lim_{m\to\infty}\gamma_m^{\mathrm{geo}}
=
\log 2-h_\partial-c_{u\to w}.
\]
两式相减即得最后结论。证毕。
\end{proof}

\begin{corollary}[双口径残差码率的同步塌缩判据]\label{cor:conclusion-double-kl-residual-rate-synchronous-collapse}
\leanverified{paper\_conclusion\_double\_kl\_residual\_rate\_synchronous\_collapse}
沿用定理 \ref{thm:conclusion-double-kl-residual-conditional-complexity-identity} 的设定，则以下两条件等价：
\[
c_{u\to w}=c_{w\to u}=0;
\]
\[
\lim_{m\to\infty}\frac1m\,\EE_\omega R_m(\omega)
=
\lim_{m\to\infty}\frac{1}{m\log 2}\sum_{x\in X_m}u_m(x)\log d_m(x)
=
\frac{\log 2-h_\partial}{\log 2}.
\]
亦即，只有当双向 KL 线性常数同时消失时，微态均匀口径的最优条件复杂度密度与边界均匀口径的理想枚举密度才会同步塌缩到同一算术基线。
\end{corollary}
\begin{proof}
由定理 \ref{thm:conclusion-double-kl-residual-conditional-complexity-identity}，两极限分别等于
\[
\frac{\log 2-h_\partial+c_{w\to u}}{\log 2},
\qquad
\frac{\log 2-h_\partial-c_{u\to w}}{\log 2}.
\]
故它们同时等于基线 \((\log 2-h_\partial)/\log 2\) 当且仅当
\[
c_{w\to u}=0,\qquad c_{u\to w}=0.
\]
证毕。
\end{proof}

\begin{theorem}[中心规范熵在全纤维对称中的零占比]\label{thm:conclusion-central-gauge-entropy-zero-share}
设 \(b(m)\) 为所有大小恰为 \(2\) 的纤维数，并记其中心规范熵损失为
\[
\Delta S_{\mathrm{gauge}}(m)=(b(m)-1)\log 2.
\]
又记全纤维对称群为 \(G(\Fold_m)\)。则有
\[
\frac{\Delta S_{\mathrm{gauge}}(m)}{\log |G(\Fold_m)|}\longrightarrow 0
\qquad (m\to\infty).
\]
\end{theorem}
\begin{proof}
由定义 \(b(m)\le |X_m|\)，而正文已给出
\[
|X_m|=F_{m+2}=\Theta(\varphi^m).
\]
故
\[
\Delta S_{\mathrm{gauge}}(m)=O(\varphi^m).
\]
另一方面，正文给出
\[
\log |G(\Fold_m)|=\Omega(m2^m).
\]
于是
\[
0\le
\frac{\Delta S_{\mathrm{gauge}}(m)}{\log |G(\Fold_m)|}
\ll
\frac{\varphi^m}{m2^m}\to 0
\qquad (\varphi<2).
\]
证毕。
\end{proof}

\begin{theorem}[Gödel 单整数化与最优残差码之间的非准等距障碍]\label{thm:conclusion-godel-single-integer-non-quasiisometric-obstruction}
\leanverified{paper\_conclusion\_godel\_single\_integer\_non\_quasiisometric\_obstruction}
设 \(\tau(\omega)\) 为正文中的事件码序列，\(G(\tau(\omega))\) 为其 Gödel 单整数化结果。若某个集合 \(E_m\subset\Omega_m\) 上成立
\[
T(\omega)\ge \alpha m
\qquad (\omega\in E_m),
\]
其中 \(\alpha>0\) 且 \(T(\omega):=|\tau(\omega)|\)，则在 \(E_m\) 上有一致下界
\[
\frac{\log_2 G(\tau(\omega))}{K(r_{m,S}(\omega)\mid X)}
\ge
c_\alpha \log_2 m-O(1).
\]
特别地，不存在把残差外置压缩成单个 Gödel 整数而仍与最优条件残差码保持准等距的统一机制。
\end{theorem}
\begin{proof}
正文给出下界
\[
\log_2 G(\tau(\omega))
\ge
T(\omega)\log_2 T(\omega)-c_0T(\omega)-c_1.
\]
而最优条件残差码满足
\[
K(r_{m,S}(\omega)\mid X)\le \lceil \log_2 d_m(X)\rceil+c.
\]
再由 \(d_m(X)\le 2^m\)，有
\[
K(r_{m,S}(\omega)\mid X)\le m+O(1).
\]
若 \(T(\omega)\ge \alpha m\)，则
\[
\log_2 G(\tau(\omega))
\ge
\alpha m\log_2 m+O(m),
\]
从而
\[
\frac{\log_2 G(\tau(\omega))}{K(r_{m,S}(\omega)\mid X)}
\ge
\frac{\alpha m\log_2 m+O(m)}{m+O(1)}
\ge
c_\alpha \log_2 m-O(1).
\]
该比值发散，因此准等距不可能。证毕。
\end{proof}

\begin{corollary}[椭圆输运条件数相对于最优账本的对数发散]\label{cor:conclusion-elliptic-transport-condition-number-log-divergence}
在定理 \ref{thm:conclusion-godel-single-integer-non-quasiisometric-obstruction} 的设定下，记面积保持椭圆输运
\[
A_{\tau}:=\diag\!\bigl(G(\tau)^{1/2},G(\tau)^{-1/2}\bigr).
\]
则在 \(E_m\) 上有
\[
\frac{\log_2 \kappa(A_{\tau(\omega)})}{K(r_{m,S}(\omega)\mid X)}
\ge
c_\alpha \log_2 m-O(1).
\]
因此，单整数椭圆输运所引入的几何畸变，其对数量级相对于最优残差账本必然发散。
\end{corollary}
\begin{proof}
正文断言
\[
\kappa(A_\tau)=G(\tau).
\]
故
\[
\log_2 \kappa(A_{\tau(\omega)})=\log_2 G(\tau(\omega)).
\]
代入定理 \ref{thm:conclusion-godel-single-integer-non-quasiisometric-obstruction} 即得。证毕。
\end{proof}

\begin{theorem}[椭圆输运的 Jordan--H{\"o}lder 二主缺陷定理]\label{thm:conclusion-elliptic-transport-jordan-holder-two-primary-defect}
设 \(P(z)\in K[z]\) 可分且 \(a_0a_n\neq 0\)，\(Q_r(w)\) 为正文定义的椭圆输运结果式。记 \(L/K\) 为 \(P\) 的分裂域。
\begin{enumerate}
  \item 若 \(r\) 超越于 \(K\)，则
  \[
  \Gal(Q_r/K(r))\cong \Gal(P/K).
  \]
  \item 若 \(r\in K^\times\) 为代数尺度，记 \(M/K\) 为 \(Q_r\) 的分裂域，则 \(\Gal(L/M)\) 为一个 \(2\)-群；因此 \(\Gal(M/K)\) 的 Jordan--H{\"o}lder 因子多重集，恰由 \(\Gal(P/K)\) 的 Jordan--H{\"o}lder 因子多重集删去若干个 \(C_2\) 后得到。特别地，一切奇阶单因子与一切非阿贝尔单因子全部保留。
\end{enumerate}
\end{theorem}
\begin{proof}
第一条正是正文已经证明的超越尺度不变性。

对第二条，正文给出夹域关系
\[
M\subseteq L\subseteq M\bigl(\sqrt{w_1^2-4},\dots,\sqrt{w_n^2-4}\bigr),
\]
故
\[
[L:M]\mid 2^n,
\]
从而 \(\Gal(L/M)\) 为 \(2\)-群。又由于 \(L/K\) 为 Galois、\(M\) 为其中间域，存在短正合列
\[
1\longrightarrow \Gal(L/M)
\longrightarrow \Gal(L/K)
\longrightarrow \Gal(M/K)
\longrightarrow 1.
\]
而 \(2\)-群的全部 Jordan--H{\"o}lder 因子都同构于 \(C_2\)。因此 \(\Gal(M/K)\) 的 Jordan--H{\"o}lder 因子恰由 \(\Gal(L/K)=\Gal(P/K)\) 的 Jordan--H{\"o}lder 因子删去若干个 \(C_2\) 所得；尤其奇阶单因子与非阿贝尔单因子不会消失。证毕。
\end{proof}

\begin{corollary}[非可解性在代数尺度下不可被二值输运洗去]\label{cor:conclusion-elliptic-transport-nonsolvability-persists}
在定理 \ref{thm:conclusion-elliptic-transport-jordan-holder-two-primary-defect} 的设定下，若 \(\Gal(P/K)\) 非可解，则对任意代数尺度 \(r\in K^\times\)，群
\[
\Gal(Q_r/K)
\]
仍非可解。
\end{corollary}
\begin{proof}
非可解有限群的 Jordan--H{\"o}lder 因子中必含至少一个非阿贝尔单群。由定理 \ref{thm:conclusion-elliptic-transport-jordan-holder-two-primary-defect}，代数尺度输运至多删去若干个 \(C_2\) 因子，而不会删去任何非阿贝尔单因子。故 \(\Gal(Q_r/K)\) 仍含非阿贝尔单因子，从而仍非可解。证毕。
\end{proof}

\begin{theorem}[分支账本的 Abel 可观测赤字]\label{thm:conclusion-branch-ledger-abelian-observable-deficit}
设
\[
W=(\ZZ/2\ZZ)^I\rtimes G
\]
为正文中的分支--对称半直积，记 \(n:=|I|\)，并设 \(G\) 在 \(I\) 上分成 \(k\) 个轨道。则
\[
W^{ab}\cong (\ZZ/2\ZZ)^k\oplus G^{ab}.
\]
尤其，原始点态分支账本含有 \(n\) 个二值坐标，而 Abel 可观测层只能保留其中 \(k\) 个；因此 Abel 可观测赤字恰为
\[
n-k=\sum_{j=1}^{k}(|O_j|-1),
\]
其中 \(O_1,\dots,O_k\) 为轨道分解。特别地，若作用是传递的，则 \(k=1\)，故 \(n-1\) 个分支位在一切 Abel 观测中都不可见，只剩总奇偶这一单比特。
\end{theorem}
\begin{proof}
正文已给出
\[
W^{ab}\cong V_G\oplus G^{ab},
\qquad
V_G\cong (\ZZ/2\ZZ)^k.
\]
故 Abel 层可见的分支部分恰有 \(k\) 个独立 \(\FF_2\)-坐标。与原始 \(n\) 个点态分支位相比，损失正是 \(n-k\)。传递情形取 \(k=1\) 即得最后结论。证毕。
\end{proof}

\begin{theorem}[有限 CMV 证书塔上的 Pad\'e--Hardy 有限冻结]\label{thm:conclusion-finite-cmv-certificate-pade-hardy-freezing}
\leanverified{paper\_conclusion\_finite\_cmv\_certificate\_pade\_hardy\_freezing}
设 \(F_\infty\) 为极限 Carath\'eodory 函数，\(F_L\) 为第 \(L\) 级有限 CMV 证书所对应的 Carath\'eodory 函数。固定非负整数 \(M,N\)，并假定 \(F_\infty\) 在 \(0\) 处的 \([M/N]\)-Pad\'e 逼近存在且唯一。则存在 \(L^\ast(M,N)\)，使得对一切 \(L\ge L^\ast(M,N)\)，\(F_L\) 在 \(0\) 处的 \([M/N]\)-Pad\'e 逼近不仅存在且唯一，而且与 \(F_\infty\) 的 \([M/N]\)-Pad\'e 逼近严格相同。

从而，对任何与极限分母零点集无交的紧集 \(K\subset\CC\)，该 Pad\'e 逼近在 \(K\) 内的零点数与极点数自某一有限级别起完全冻结。
\end{theorem}
\begin{proof}
\([M/N]\)-Pad\'e 逼近的系数是 \(F\) 在 \(0\) 处前 \(M+N+1\) 个 Taylor 系数的有理函数；其存在唯一性由相应 Hankel 行列式非零刻画。

而正文的有限 CMV 证书冻结链表明：任意有限个 Taylor 系数在某一有限级别 \(L^\ast\) 之后完全稳定，即
\[
[F_L]_{0}^{(j)}=[F_\infty]_{0}^{(j)}
\qquad (0\le j\le M+N).
\]
因此 Pad\'e 构造所需的全部代数数据在 \(L\ge L^\ast\) 后逐项一致，故 \([M/N]\)-Pad\'e 逼近本身也完全一致。其在固定紧集 \(K\) 上的零点数与极点数作为同一有理函数的局部数据，自然随之冻结。证毕。
\end{proof}

\begin{corollary}[局部有理谱谓词的有限可判定性]\label{cor:conclusion-local-rational-spectral-predicates-finite-decidability}
\leanverified{paper\_conclusion\_local\_rational\_spectral\_predicates\_finite\_decidability}
在定理 \ref{thm:conclusion-finite-cmv-certificate-pade-hardy-freezing} 的设定下，任意仅依赖于有限阶局部谱矩、或等价地仅依赖于 Carath\'eodory 函数在 \(0\) 处有限个 Taylor 系数的布尔谓词，沿有限 CMV 证书塔都在某一有限级别之后稳定为常值。

特别地，凡可表述为有限个 Pad\'e 分母非退化条件、有限个零点/极点位置不等式、或有限个局部谱矩多项式关系的命题，均可在有限级别上完成离线判定，而无须取极限对象。
\end{corollary}
\begin{proof}
由正文的 CMV 冻结定理，任意有限多个局部谱矩或 Taylor 系数都在某一有限级别之后逐项稳定。任何此类布尔谓词都由有限维代数运算与有限个等式/不等式组合而成，故其真值在同一有限级别之后不再改变。证毕。
\end{proof}

\noindent
因此，残差外置、Gödel 单整数化、椭圆输运与有限 CMV 证书在本文对象中并非彼此独立的四套机制：双向 KL 常数把最优条件复杂度密度与理想枚举码率精确焊接起来；Gödel 单整数化与面积保持椭圆输运则共同暴露出相对于最优残差账本不可压缩的对数鸿沟；在代数尺度下，输运只能剥离若干个 \(C_2\) 因子，却不能洗去非阿贝尔或奇数部分的 Galois 复杂性；而一旦进入有限 CMV 证书塔的点态稳定化区间，一切有限阶 Pad\'e--Hardy 局部对象又都会在某一有限级别后冻结为极限值。于是，信息论、Galois 结构与局部谱几何在这里被统一压缩为一条同时具备外置下界、代数刚性与有限可判定性的终端链。

\endinput
